\section{CST Model Construction for the \texorpdfstring{\(2\times1\)}{2x1} Patch Array}
\label{sec:cst-array-model}

The \(2\times1\) array model was constructed from the tuned single-patch element.  The final
single-patch geometry was duplicated in the \(x\)-direction, and the center-to-center spacing was set
to
\[
 d = 59.935~\mathrm{mm}\approx \frac{\lambda_0}{2}.
\]
The substrate and ground plane were extended so that both elements and both feed lines were fully
supported.  Each element was fed by its own waveguide port so that the two ports could be excited
independently for matching, coupling, and beam-steering analysis.

The same mesh-control simplification used for the final single-patch run was retained: the metal
regions were modeled as PEC, while the DiClad 880-IM dielectric loss was retained.  This approach
kept the CST model below the mesh-cell limitation while preserving the dominant resonant and
radiation behavior needed for the antenna and array analysis.

\subsection{Final Array Geometry Parameters}

\begin{table}[H]
\centering
\caption{Final tuned CST dimensions for the \(2\times1\) array model.}
\label{tab:cst-array-final-parameters}
\begin{tabular}{@{}ll@{}}
\toprule
\textbf{Parameter} & \textbf{Final Value} \\
\midrule
Target frequency, \(f_0\) & \(2.501~\mathrm{GHz}\) \\
Substrate & DiClad 880-IM \\
Substrate thickness, \(h\) & \(1.524~\mathrm{mm}\) \\
Patch width, \texttt{patchW} & \(47.382~\mathrm{mm}\) \\
Final array patch length, \texttt{patchL} & \(36.70~\mathrm{mm}\) \\
Feed width, \texttt{feedW} & \(4.735~\mathrm{mm}\) \\
Feed length, \texttt{feedL} & \(15.000~\mathrm{mm}\) \\
Inset depth, \texttt{insetY} & \(13.35~\mathrm{mm}\) \\
Inset gap, \texttt{insetGap} & \(1.000~\mathrm{mm}\) \\
Element spacing, \texttt{arrayD} & \(59.935~\mathrm{mm}\) \\
Metallization model for tuned run & PEC approximation \\
\bottomrule
\end{tabular}
\end{table}

\subsection{CST Array Construction Notes}

The first array run used the final tuned single-patch dimensions directly.  The array resonance then
shifted downward because the second element changed the electromagnetic environment of the first
element.  This behavior is expected in a closely spaced array because the input impedance of each
radiating element is affected by mutual coupling.  Therefore, the array required a separate tuning
sequence after the single-patch tuning was completed.

The final CST array model used two waveguide ports.  Port~1 excited the first element and Port~2
excited the duplicated element.  The two-port model allowed CST to compute \(S_{11}\), \(S_{22}\),
\(S_{12}\), and \(S_{21}\), which were used to document input matching and mutual coupling.
